%\documentclass[notes,10pt,aspectratio=169]{beamer}
\documentclass[10pt,aspectratio=169]{beamer}

\usetheme[progressbar=frametitle]{metropolis}
\usecolortheme{rose}
\usecolortheme{default}

\usepackage[utf8]{inputenc}
\usepackage[T1]{fontenc}
\usepackage{lmodern}
\usepackage{xcolor}
\usepackage{tikz}
\usepackage{booktabs}
\usetikzlibrary{shapes.geometric, arrows, positioning}

\tikzstyle{block} = [rectangle, draw, text width=4cm, align=center, rounded corners, minimum height=1cm]
\tikzstyle{decision} = [rectangle, draw, text width=5cm, align=center, fill=blue!10, rounded corners, minimum height=1cm]
\tikzstyle{terminal} = [rectangle, draw, text width=4.5cm, align=center, fill=yellow!30, rounded corners, minimum height=1cm]
\tikzstyle{end} = [rectangle, draw, text width=5cm, align=center, fill=green!30, rounded corners, minimum height=1cm]
\tikzstyle{arrow} = [->, thick]

\usepackage{adjustbox}
\setbeamertemplate{itemize item}[triangle]

\definecolor{accent}{RGB}{78,205,196}
\setbeamercolor{frametitle}{fg=black,bg=white}
\setbeamercolor{item}{fg= orange!80}
\setbeamercolor{button}{bg=orange, fg=white}

%% BibLaTeX disabled (no citations in this presentation)
% \usepackage[backend=biber, style=apa, citestyle=authoryear]{biblatex}
% \addbibresource{../references.bib}

\title{Meeting with Phil}
\author{Lucas Schmitz\inst{1}}
\institute{\inst{1} Yale University}
\date{April 2, 2026}

\begin{document}


%-------------------------------
\begin{frame}{Last Meeting (6 February)}
\textbf{What we discussed:}
\begin{itemize}
    \item Talked about cross-selling  
    \item Presented a model of switching costs with persistence in preferences
\end{itemize}

\vspace{0.3cm}
\textbf{Today's meeting:}
\begin{itemize}
    \item Economics of policies that change lenders' information sets
    \item Data we would have access to
    \item Structural model for estimation
\end{itemize}
\end{frame}

%-------------------------------
\begin{frame}{Policies that affect the amount of public information for lenders}


\textbf{Internationally:}
\begin{itemize}
    \item Fair Credit Reporting Act (FCRA) sets  Time limits (e.g. 10 years for bankruptcies)
    \item Open Banking Policies (e.g. in the UK) require banks to share customer data (including balances and transactions) with third parties  
\end{itemize}

\textbf{In Chile:}
\begin{itemize}
    \item Restrictions to credit bureaus information. They deleted delinquencies for short term unemployed individuals, excluded credit inquiries from credit scores and once they erased default histories for amounts below \$1,500. 
    \item Expand the reporting entities: incorporated credit cards by retailers (law 21,130) and incorporated other non-bank lenders (law 21,680)
\end{itemize}
\end{frame}


%-------------------------------
%-------------------------------
\begin{frame}{Economics of changing lenders' information sets }
\begin{itemize}
    \item \textbf{Setup:} Incumbent bank has private info on borrowers, rival banks do not
    \item Effects of policies that expand public information:
    \begin{enumerate}
        \item \textbf{Reduces adverse selection:} switchers are no longer assumed to be lemons 
        \item \textbf{Reduces incumbent's information rents:} Incumbent's advantage came from knowing $h$ (default probability) when rivals did not. If rivals can also screen, the ``winner's curse'' weakens $\Rightarrow$ incumbent's market power falls
        \item \textbf{Affects negatively risky borrowers:} uninformed lenders may raise their rates or deny them credit
    \end{enumerate}
\end{itemize}
\textbf{Predictions:}
\begin{itemize}
    \item[P1.] Switching rates increase for good borrowers (rival banks price more aggressively)
    \item[P2.] Interest rates fall for good borrowers (separating $\to$ competitive pricing)
    \item[P3.] Default rates of switchers decrease  (reduction of adverse selection) 
    \item[P4.] Negative effect on risky borrowers: may face higher rates or lose access
\end{itemize}
\end{frame}




%-------------------------------
%-------------------------------
\begin{frame}{Data: Chilean Administrative Credit Records (CMF)}

\vspace{0.3cm}
\textbf{Individual-level data on:}
\begin{itemize}
    \item \textbf{Credit cards} (bank \& non-bank): usage, payments, delinquency
    \item \textbf{Loans:} amounts, interest rates, payment plans, maturities, delinquency
    \item \textbf{Accounts} (checking, savings): balances and transactions
\end{itemize}

\textbf{Unit of observation:} borrower $\times$ lender $\times$ month.
\end{frame}


%-------------------------------
%-------------------------------
\begin{frame}{Structural Model: Two-Period Game}
 
\textbf{Period 1 --- First loan:}
\begin{itemize}
    \item Two banks compete on rate $r_{j1}$; neither observes borrower type $h \sim F(h)$
    \item Winning bank learns $h$ from the relationship; borrower develops switching cost $\lambda$
\end{itemize}

\vspace{0.2cm}
\textbf{Period 2 --- Second loan:}
\begin{itemize}
    \item Incumbent  sets rate $r_{12} = \sigma^*(h)$ and rival sets $r_{22} = r_2^*$
    \item Borrower chooses via demand: $\; s(r_{12}, r_{22}) = \dfrac{1}{1 + \exp(\alpha(r_{12} - r_{22}) - \lambda)}$
\end{itemize}

\vspace{0.2cm}
\textbf{Backward induction:} Period-2 equilibrium yields continuation values $V^I > V^R > 0$. In Period 1, each bank solves:
$$\max_{r_{j1}} \;\; s(r_{j1}, r_{-j1}) \cdot \Big[\,E_h[(1-h)\,r_{j1} - 1] + \Delta V\,\Big], \qquad \Delta V = V^I - V^R$$
Banks price Period-1 loans below break-even to secure the incumbent position in Period 2.
\end{frame}


%-------------------------------
%-------------------------------
\begin{frame}{Structural Model: Period-2 Equilibrium}

\textbf{Incumbent (bank 1)} prices each type at break-even $+$ markup:
$$\sigma^*(h) = \underbrace{\frac{1}{1-h}}_{\text{break-even}} + \underbrace{\frac{1}{\alpha\big(1-s(\sigma^*(h), r_2^*)\big)}}_{\text{markup}}$$

\vspace{0.2cm}
\textbf{Rival (bank 2)} sets a single rate:
$$r_2^* = \frac{1}{1 - \tilde{h}} + \frac{1}{\alpha}\cdot\frac{\int (1-s)(1-h)\, dF}{\int s\,(1-s)(1-h)\, dF}$$

\vspace{0.2cm}
\textbf{Primitives to estimate:} $\alpha$ (price sensitivity), $\lambda$ (switching cost), $F(h)$ (type distribution)
\end{frame}


%-------------------------------
%-------------------------------
\begin{frame}{Structural Model: Identification}
\begin{itemize}
    \item Observables: winning rate $r_i^{\text{win}}$, winning bank $W_i \in \{1,2\}$, default $D_i \in \{0,1\}$
    \item Sequential identification:
    \begin{enumerate}
        \item \textbf{Rival rate:} $r_2^*$ identified as the common rate among bank-2 winners
        \item \textbf{Type recovery:} $E[D_i \mid r_i^{\text{win}} = r,\, W_i = 1] = h$ \; $\Rightarrow$ invert $\sigma^*$ to recover each borrower's type    
        \item \textbf{$\alpha$ and $\lambda$:} From the equilibrium FOC, for any two types $h_a \neq h_b$:
        $$\ln\!\big(\alpha\, m(h) - 1\big) \;=\; \lambda \;-\; \alpha\,\Delta(h)$$
        where $m(h) = \sigma^*(h) - \tfrac{1}{1-h}$ is the markup and $\Delta(h) = \sigma^*(h) - r_2^*$ is the rate gap. Slope identifies $\alpha$; intercept identifies $\lambda$.
    \item \textbf{Type distribution:} $F(h)$ recovered from bank-1 winners, correcting for selection via $s(h)$
\end{enumerate}
\end{itemize}
\end{frame}


 
\end{document}
